In this thesis we provided the first constant-time implementation of HQC and the
first masked implementation of HQC's public key encryption scheme.

To do so, we evaluated different multiplication methods in terms of both performance and constant-time behavior, 
crowning the Frobenius Additive FFT as the best one on the ARM Cortex-M4 and the Karatsuba algorithm on x86 architectures;
moreover, we implemented a masking scheme for the key generation, encryption and decryption primitives
of HQC's public key encryption scheme.

We provided a comparative analysis of the Round 4 NIST PQC candidates, showing how both our constant-time and masked implementations 
of HQC offer a good trade-off in terms of memory consumption and execution time, making them suitable for the standardization.

As future work, it could be possible to run more accurate tests in order to evaluate the efficacy of our masking scheme against power attacks.
Since our masking scheme targets the PKE only, it could also be possible to extend the masking scheme to 
the SHAKE functions and make the whole KEM secure against power attacks, such as the one proposed in~\cite{DBLP:journals/tches/UenoXTITH22}.
